\chapter{Echocardiography}
\section*{What Is This Test?} Echocardiography uses ultrasound to produce moving images of the heart, valves, chambers, muscle, and blood flow.
\section*{Alternative Names} Echocardiogram; cardiac ultrasound; transthoracic echo; TTE.
\section*{Principle} Two-dimensional, Doppler, and sometimes three-dimensional ultrasound assess structure, motion, pressure estimates, and flow.
\section*{Why This Test Is Ordered} It evaluates murmurs, breathlessness, heart failure, valve disease, cardiomyopathy, congenital disease, clots, infection, or pulmonary-pressure concerns.
\section*{Primary vs Secondary Test} It is a primary noninvasive cardiac imaging test; transesophageal echo, CT, MRI, or catheterization may answer additional questions.
\section*{How This Test Is Performed} A transducer is moved over the chest. A transesophageal study uses a sedated probe in the esophagus.
\section*{What Preparation Is Needed} Transthoracic testing usually needs none; transesophageal testing requires fasting and an escort after sedation.
\section*{Understanding Results} Reports describe chamber size, pumping function, wall motion, valves, fluid, and flow. Findings require clinical correlation.
\section*{Follow-up Tests} ECG, stress testing, cardiac MRI or CT, repeat echocardiography, or cardiology consultation may follow.
\section*{References} \begin{enumerate}\item MedlinePlus. Echocardiography.\item American Society of Echocardiography. Echocardiography standards.\end{enumerate}
\clearpage\section*{Understanding Results and Follow-up}
The laboratory report should identify the specimen, method, units, reference interval or decision limit, and whether the result is positive, negative, abnormal, or inconclusive. Reference intervals vary with age, sex, pregnancy, specimen type, assay, and laboratory; a result outside a range is not automatically a diagnosis, and a result within a range does not always exclude disease. Discuss unexpected results with the ordering clinician, who can decide whether repeat or confirmatory testing, treatment, or referral is appropriate. Do not change medicines or delay urgent care based on an isolated result.
\section*{Regulatory Status and Authorization Date}This chapter describes a test category, not a specific commercial product. FDA, CDSCO, EU IVDR, and MHRA approval or registration dates therefore cannot be assigned without the assay manufacturer, product name, instrument, and intended use. For laboratory-developed tests, an individual agency approval date may not exist. See the Preface for the interpretation of regulatory status.
\clearpage
